\documentclass[11pt,a4paper]{article}
\usepackage[T1]{fontenc}
\usepackage[utf8]{inputenc}
\usepackage[french]{babel}
\usepackage{lmodern}
\usepackage{amsmath,amssymb}
\usepackage[margin=2.35cm]{geometry}
\usepackage{booktabs,longtable,tabularx}
\usepackage{enumitem}
\usepackage{xcolor}
\usepackage{microtype}
\usepackage{hyperref}

\hypersetup{colorlinks=true,linkcolor=blue!45!black,urlcolor=blue!45!black}
\setlist{nosep}
\newcommand{\code}[1]{\texttt{#1}}
\providecommand{\degree}{\ensuremath{^\circ}}
\sloppy

\title{Prompt M4.3Live-v2\\
\large Libérer le sens du prêt, et chercher ce qui détermine la rotation
du crédit}
\author{Lignée M4.3Live-v2 --- succède à M4.3Live
(\code{m4\_3live\_credit\_soc}, programme terminé et publié),\\
destiné à une instance Claude Opus 5}
\date{21 août 2026 --- révision 1}

\begin{document}
\maketitle
\tableofcontents
\bigskip

Destinataire : une instance Claude Opus 5 (Claude Code), à qui ce document est confié comme spécification de tâche complète et autonome. Ce n'est pas un brouillon à discuter avant de commencer : les décisions de fond sont prises ci-dessous ; les points explicitement laissés ouverts (§12) sont à trancher et documenter toi-même, pas à retourner à l'utilisateur avant d'avoir essayé.

Version : 21 août 2026, révision 1. Lignée : succède à \textbf{M4.3Live} (\code{m4\_3live\_credit\_soc/}), qui est un programme \textbf{terminé et publié} --- deux rapports PDF, 203 runs enregistrés, parité bit à bit avec M4.3 vérifiée sur 8000 pas. Convention de dossier : \code{m4\_3live\_v2\_credit\_soc/}.

\textbf{Provenance.} Ce document dérive de \code{m4\_3live\_v2\_credit\_soc/ROADMAP.md}, écrite à partir de 28 notes manuscrites portées par l'utilisateur sur \code{m4\_3live\_credit\_soc/report/rapport\_final.pdf} le 18 août 2026, puis complétée par quatre retours oraux du 21 août. La feuille de route explique le \emph{pourquoi} de chaque chantier et cite sa note d'origine ; \textbf{ce prompt donne le \emph{quoi} et l'\emph{ordre}}. Les deux se lisent ensemble : en cas de contradiction, ce prompt fait foi, et la contradiction est à signaler.

\bigskip

\setcounter{section}{-1}
\section{Lectures préalables obligatoires, et pièges du dépôt}

Avant d'écrire une ligne de code, lire dans cet ordre :

\begin{itemize}
\item \textbf{\code{m4\_3live\_v2\_credit\_soc/ROADMAP.md}} --- la feuille de route dont ce prompt est l'exécution. Elle donne, chantier par chantier, la note d'utilisateur qui le motive et les conséquences à instrumenter.
\item \textbf{\code{m4\_3live\_credit\_soc/m4\_3live/model.py}} (\textasciitilde{}1000 lignes) --- le moteur de référence. C'est le fichier dont ce prompt cite les numéros de ligne ; si le fichier a changé depuis, revérifier les citations, ne pas les croire sur parole.
\item \code{m4\_3live\_credit\_soc/m4\_3live/kernel.py} --- le noyau d'institution (table de Hermite d'ordre 4, \ensuremath{\delta}* = h(C) \ensuremath{-} K\_b). Validé, coût 1,3 \% du pas : à reprendre tel quel.
\item \code{m4\_3live\_credit\_soc/m4\_3live/live.py} --- session pilotable, journal d'interventions, rejeu, snapshots. Testé, non remis en cause.
\item \textbf{\code{m4\_3live\_credit\_soc/report/rapport\_final.pdf}} (34 pages) --- les résultats de v1. Lire au moins §5 (ablation K0), §6 (la tension et ses limites) et §7 (les deux groupes de covariance) : v2 en hérite comme acquis, et une bonne partie de son programme consiste à les prolonger.
\item \textbf{\code{m4\_3live\_credit\_soc/report/conception\_m4\_3live.pdf}} (27 pages) --- les décisions d'architecture de v1 et leurs justifications.
\item \code{m4\_3live\_credit\_soc/JOURNAL.md} --- le journal de travail de v1, y compris les pièges rencontrés et les inférences réfutées en cours de route.
\item \code{simulation\_lab/contracts.py}, \code{simulation\_lab/jobs.py}, \code{simulation\_lab/web/app.py} --- l'orchestration existante, dont v2 étend l'usage sans le casser.
\item \code{CLAUDE.md} (racine du dépôt) --- règles de travail du dépôt.
\end{itemize}

\subsection*{Deux passages de \code{CLAUDE.md} sont obsolètes et vont t'induire en erreur}

\begin{enumerate}
\item \code{CLAUDE.md} désigne \code{m4b\_credit\_soc\_mini/} comme « moteur actif ». \textbf{Faux pour ce travail} : le moteur de référence est \code{m4\_3live\_credit\_soc/m4\_3live/model.py}.
\item \code{CLAUDE.md} donne \code{random.Random(seed)} comme convention RNG. \textbf{Faux ici} : la lignée M4.3/M4.3Live utilise \code{numpy.random.default\_rng(seed)} (\code{m4\_3live/model.py:718}, \code{Simulation.\_\_init\_\_}). Reproduire cette convention dans le fork.
\end{enumerate}

\subsection*{Règle de confiance à 95 \%}

Si tu n'es pas certain à 95 \% d'un fait sur ce projet, le signaler explicitement, en distinguant \textbf{fait observé} (lu dans le code ou les données), \textbf{inférence} (déduit), \textbf{hypothèse} et \textbf{incertitude}. Citer systématiquement \code{fichier.py:ligne}. Ne jamais inventer un comportement du moteur non vérifié dans le code. Cette règle a une traduction typographique obligatoire dans les rapports (§10).

\subsection*{La règle du fork : \code{m4\_3live/} est en lecture seule, définitivement}

Le paquet \code{m4\_3live\_credit\_soc/m4\_3live/} \textbf{ne doit pas être modifié}, ni maintenant ni plus tard. Ce n'est pas une précaution de style : c'est le moteur qui a produit 203 runs enregistrés et deux rapports publiés, et toute modification invaliderait rétroactivement des résultats cités. L'utilisateur a énoncé cette contrainte explicitement, en majuscules, pendant le travail de v1.

v2 \textbf{forke} ce paquet dans son propre dossier (\code{m4\_3live\_v2\_credit\_soc/m4\_3live\_v2/}, nom à confirmer §12), par copie, pas par import. Le fork peut ensuite être modifié librement. Toute mesure nouvelle qu'on voudrait faire sur v1 doit rester un dérivé des séries déjà écrites (voir \code{m4\_3live\_credit\_soc/scripts/tension.py}, qui reconstruit la tension a posteriori sans toucher au moteur).

\subsection*{Un trou connu dans la feuille de route}

\textbf{\code{ROADMAP.md} §9 est un emplacement vide}, réservé aux annotations que l'utilisateur a portées sur \code{conception\_m4\_3live.pdf} et que la visionneuse n'a pas enregistrées (0 objet \code{/Subtype/Text} sur le fichier, vérifié par \code{pdftk dump\_data\_annots} et par \code{mutool show \dots{} grep}). \textbf{N'invente rien pour cette section.} Si l'utilisateur fournit une copie annotée en cours de travail, l'intégrer ; sinon, laisser le trou et le signaler dans le rapport de conception.

\subsection*{Citer les sections, jamais les numéros de figure}

Quand tu renvoies au rapport v1, cite \textbf{§4}, \textbf{§6.3}, \textbf{§7.1} --- jamais « figure 12 ». Les numéros de figure ont déjà bougé une fois (l'insertion d'une section en §6 a décalé toutes les figures à partir de la quinzième) et rebougeront. Les numéros de section utilisés dans ce prompt ont été vérifiés sur le PDF construit le 21 août 2026 :

\begin{center}
\begin{tabularx}{\linewidth}{@{}lX@{}}
\toprule
\textbf{§} & \textbf{Contenu} \\
\midrule
1 & Le modèle et les notations \\
2 & Ce qui est mesuré \\
3 & Protocole (3.1 fenêtre/graines, 3.2 bras, 3.3 appariement, 3.4 validation) \\
4 & Résultats de la campagne \\
5 & Ablation sur le capital de naissance K0 \\
6 & La tension (6.1 définition, 6.2 niveaux, 6.3 mortalité, 6.4 le test qui échoue, 6.5 exposant dlnT/dlnA) \\
7 & La loi d'échelle (7.1 covariance d'échelle, 7.2 trois \ensuremath{\gamma}, 7.3 covariance de pas de temps) \\
8 & Verdict --- 9 Limites --- 10 Reproduire \\
\bottomrule
\end{tabularx}
\end{center}

\bigskip

\setcounter{section}{0}
\section{Le mandat, en une phrase}

v2 porte \textbf{un changement de modèle} et \textbf{une question théorique}. Ils ne sont pas de même rang :

\begin{quote}
\textbf{Le changement de modèle est le livrable primaire ; la question théorique est le programme scientifique qu'il rend possible.}
\end{quote}

Concrètement : si le budget se resserre, livrer le §3.1 (sens du prêt libéré) mesuré proprement, avec son rapport, plutôt qu'un demi-chantier de modèle et une demi-réponse théorique. Ne pas confondre les deux dans un même verdict.

\textbf{Le changement de modèle} (§3.1) supprime la règle « dans chaque paire, la plus riche prête ». Son intérêt n'est pas cosmétique : il ouvre un régime que v1 \textbf{ne pouvait pas produire} --- une entité riche à mauvaise technologie qui cède son capital à une entité pauvre à meilleur rendement marginal et vit de l'intérêt. C'est ce régime qu'il faut caractériser.

\textbf{La question théorique} (§5) : \emph{quelle combinaison de (A, \ensuremath{\gamma}, \ensuremath{\delta}, K0, \ensuremath{\lambda}, \ensuremath{\rho}) détermine la rotation du crédit \code{loan\_volume/K\_tot} à l'état stationnaire ?} v1 a montré que cette rotation prédit la mortalité à 0,4 \% près quel que soit le levier employé (R\ensuremath{^2} = 0,998 sur quatre leviers indépendants), là où la « tension » échoue. Savoir la prédire donnerait la chaîne rotation \ensuremath{\rightarrow} mortalité \ensuremath{\rightarrow} population \ensuremath{\rightarrow} production.

\bigskip

\setcounter{section}{1}
\section{Ce que v1 a établi : acquis à reprendre, pas à refaire}

Ces résultats sont \textbf{publiés et vérifiés}. Les citer, ne pas les reproduire.

\begin{center}
\begin{tabularx}{\linewidth}{@{}XXX@{}}
\toprule
\textbf{Acquis} & \textbf{Où} & \textbf{Valeur} \\
\midrule
Parité bit à bit avec M4.3, 8000 pas \ensuremath{\times} 26 colonnes & \code{tests/test\_parity\_m4\_3.py -{}-full} & écart maximal \textbf{nul}, 9 366 225 appels au noyau, 1077 s \\
Noyau d'institution (\ensuremath{\delta}* = h(C) \ensuremath{-} K\_b, Hermite ordre 4) & \code{m4\_3live/kernel.py} & validé, 1,3 \% du coût d'un pas \\
Verdict rebond & rapport §4, §8 & portée partielle court horizon : rebond, \ensuremath{\varepsilon} \ensuremath{\approx} 2,3--2,9 ; portée globale à K0 fixe : effet \textbf{inverse}, \ensuremath{\varepsilon} = 0,76 ; K0 compensé : \ensuremath{\varepsilon} = 2,02 \\
Le signe du verdict global dépend entièrement de K0 & rapport §5 & à K0 fixe, A\ensuremath{\times}1,5 \textbf{contracte} la population de 25 \% ; compenser K0 par A\textasciicircum{}\{1/(1\ensuremath{-}\ensuremath{\gamma})\} annule tout \\
Covariance d'échelle démontrée & rapport §7.1 & A\ensuremath{\rightarrow}A\ensuremath{'}, K0\ensuremath{\rightarrow}cK0 avec c = (A\ensuremath{'}/A)\textasciicircum{}\{1/(1\ensuremath{-}\ensuremath{\gamma})\} donne la \textbf{même simulation} à l'échelle c ; écart 5\ensuremath{\cdot}10\ensuremath{^-}\ensuremath{^1}\ensuremath{^5}, population rigoureusement identique \\
Covariance de pas de temps & rapport §7.3 & \ensuremath{\lambda}\ensuremath{\rightarrow}s\ensuremath{\lambda}, \ensuremath{\delta}\ensuremath{\rightarrow}1\ensuremath{-}(1\ensuremath{-}\ensuremath{\delta})\textasciicircum{}s, \ensuremath{\sigma}\ensuremath{\rightarrow}\ensuremath{\sigma}\ensuremath{\sqrt{\;}}s, \textbf{A\ensuremath{\rightarrow}sA, \ensuremath{\rho}\ensuremath{\rightarrow}s\ensuremath{\rho}} ; résidu 6--17 \% dû à la phase de marché, pas à la discrétisation \\
La tension T = K\_aut/K\_eq n'est \textbf{pas} un paramètre d'état & rapport §6.4 & excellente à \ensuremath{\delta} fixé (R\ensuremath{^2} = 0,989 sur \ensuremath{\times}15) ; mise en échec par \ensuremath{\delta} (facteur 8 sur \ensuremath{\delta} \ensuremath{\rightarrow} facteur 13 sur T, 14 \% sur la mortalité) \\
La rotation du crédit aligne les quatre leviers & rapport §6.4 & \code{morts/pop \ensuremath{\propto} (loan\_volume/K\_tot)\textasciicircum{}1,337}, R\ensuremath{^2} = 0,9982, biais par levier entre \ensuremath{-}1,7 \% et +0,7 \% \\
L'exposant dlnT/dlnA & rapport §6.5 & 0,776 / 1,145 / 1,673 pour \ensuremath{\gamma} = 0,4/0,5/0,6 ; = 0 sous compensation de K0 \\
203 runs consultables & \code{simulation\_lab}, lignée \code{m4\_3live\_credit\_soc} & annexe de traçabilité à 207 lignes, 0 sans rôle \\
\bottomrule
\end{tabularx}
\end{center}

\textbf{Définitions employées ci-dessus, à reprendre telles quelles} (rapport §6.1) :

\begin{itemize}
\item \textbf{Échelle autarcique} \code{K\_aut = [A(1\ensuremath{-}\ensuremath{\delta})/\ensuremath{\delta}]\textasciicircum{}\{1/(1\ensuremath{-}\ensuremath{\gamma})\}} : le point fixe de \code{(K + AK\textasciicircum{}\ensuremath{\gamma})(1\ensuremath{-}\ensuremath{\delta}) = K}, c'est-à-dire le capital qu'atteindrait une entité seule, sans crédit ni faillite.
\item \textbf{Capital équivalent} \code{K\_eq = (prod/(n\ensuremath{\cdot}A))\textasciicircum{}\{1/\ensuremath{\gamma}\}} : le capital qui reproduirait la production observée si toutes les entités vivantes étaient identiques.
\item \textbf{Tension} \code{T = K\_aut / K\_eq} : de combien le système tourne en dessous de son échelle autarcique. Vaut \ensuremath{\approx} 13,3 au régime de référence.
\item \textbf{Rotation du crédit} \code{loan\_volume / K\_tot} : le principal transféré au pas, rapporté au stock de capital. Variable \textbf{endogène} --- c'est une régularité descriptive, pas une loi de contrôle, et v2 doit le rappeler chaque fois qu'il l'emploie.
\end{itemize}

\bigskip

\setcounter{section}{2}
\section{Les quatre changements de modèle}

\subsection*{3.1 Le sens du prêt cesse d'être imposé --- \textbf{le chantier principal}}

\textbf{État v1.} Dans chaque paire tirée, la plus riche est prêteuse et la plus pauvre emprunteuse (\code{m4\_3live/model.py:515-520}). Quand l'optimum de production jointe voudrait faire circuler le capital du pauvre vers le riche (\ensuremath{\delta}* \ensuremath{\leq} 0), la paire est \textbf{refusée} et comptée dans \code{mkt\_blocked\_dir} (\code{model.py:525-530}). Ce compteur se stabilise autour de 25 \% des rondes.

\textbf{Ce que l'utilisateur en dit} (notes [4] et [6]) : « \emph{C'est dommage. Si l'optimisation convexe donne un sens, alors c'est celui-ci qui devrait être utilisé. Cette règle n'a pas de sens.} » et « \emph{Abandon de la nomenclature K\_\ensuremath{\ell} K\_b pour lender/borrower, passage à K\_a, K\_b.} »

\textbf{Changement.} La paire devient non ordonnée \code{(a, b)}, de capitaux \code{K\_a, K\_b}, de capital joint \code{C = K\_a + K\_b}. Le noyau retourne l'allocation optimale \code{K\_a* = h(C)} ; le transfert est \code{\ensuremath{\delta} = K\_a* \ensuremath{-} K\_a}, \textbf{de signe libre}. Le prêt va de qui cède vers qui reçoit ; le receveur est le débiteur.

\textbf{Nomenclature.} Renommer partout \code{lender}/\code{borrower} en \code{a}/\code{b}. Ce n'est pas cosmétique : tant que le code dit « prêteuse », un lecteur suppose la règle de richesse, et une régression s'y cachera.

\textbf{Conséquences à instrumenter, pas à supposer.}

\begin{itemize}
\item \code{mkt\_blocked\_dir} tombe à zéro par construction. \textbf{Garder le compteur} et l'alimenter avec ce qu'il aurait valu sous l'ancienne règle --- un \textbf{compteur contrefactuel}. Il mesure directement combien d'échanges v1 s'interdisait, et c'est une des mesures les plus parlantes du chantier.
\item Le volume de prêt augmente mécaniquement. Toute comparaison v1/v2 sur \code{loan\_volume} doit se faire à protocole apparié, jamais en absolu.
\item Le régime nouveau (riche à mauvaise technologie qui devient créancière) demande ses propres diagnostics : \textbf{part du capital détenue par les créancières nettes}, et \textbf{corrélation entre technologie et position nette}. Les ajouter aux séries.
\item Le taux d'intérêt (\code{pair\_rate}, moyenne géométrique des rendements marginaux \ensuremath{\gamma}AK\textasciicircum{}\{\ensuremath{\gamma}\ensuremath{-}1\} de chaque côté) a-t-il encore le même sens quand la créancière est la plus fragile des deux ? Question ouverte (§12), à trancher et documenter.
\end{itemize}

\textbf{Parité : hypothèse à retester, pas acquise.} En régime homogène, l'optimum est le partage égal, donc \code{\ensuremath{\delta} = (K\_b \ensuremath{-} K\_a)/2} : magnitude et sens coïncident avec l'ancienne règle. La parité bit à bit \emph{devrait} survivre. C'est une \textbf{hypothèse} : l'ordre de sommation et l'ordre d'insertion dans le carnet dépendent de l'ordre du couple, et c'est exactement le genre de détail qui décale un flottant. \code{test\_parity\_m4\_3.py -{}-full} est le premier test à repasser, et \textbf{un écart doit être expliqué ligne à ligne, pas absorbé dans une tolérance}.

\subsection*{3.2 Ordre des phases : dépréciation avant service des intérêts, en option}

\textbf{État v1.} Ordre du pas : naissances \ensuremath{\rightarrow} choc \ensuremath{\rightarrow} production \ensuremath{\rightarrow} \textbf{service des intérêts} \ensuremath{\rightarrow} \textbf{dépréciation} \ensuremath{\rightarrow} marché \ensuremath{\rightarrow} faillites (\code{m4\_3live/model.py:864-934}).

\textbf{Note [3]} : « \emph{Test à faire, échanger 5 et 6, ça devrait faire augmenter le taux de défaut.} »

\textbf{Changement.} Un champ de configuration \code{phase\_order \ensuremath{\in} \{"v1", "deprec\_first"\}}, valeur par défaut \code{"v1"} --- la parité doit rester atteignable sans détour. Sous \code{"deprec\_first"} : dépréciation puis service.

\textbf{Prédiction à falsifier, à écrire avant de lancer.} Le capital disponible au moment de payer est réduit d'un facteur (1\ensuremath{-}\ensuremath{\delta}) ; toute emprunteuse dont le capital était compris entre \code{dû} et \code{dû/(1\ensuremath{-}\ensuremath{\delta})} bascule en défaut de liquidité. À \ensuremath{\delta} = 0,01 c'est une fenêtre étroite : l'effet attendu est \textbf{petit et calculable a priori} à partir de la distribution du ratio capital/dû, qu'il faut donc enregistrer avant de mesurer. Test A/B apparié par graine.

\subsection*{3.3 Aucun traitement du cas partiel --- \textbf{périmètre retiré}}

Consigne explicite de l'utilisateur, 21 août 2026 : « \emph{Pour la v2, on se souvient, je ne veux aucun traitement du cas partiel.} »

\textbf{Ce qui sort du périmètre} : la portée \code{fraction} et toute campagne qui l'emploie ; les bras \code{null} appariés qui lui servaient de référence ; l'inversion part ex post \ensuremath{\rightarrow} part ex ante ; la question « le long horizon est-il tranchable ? » posée en portée partielle. \textbf{v2 mesure en portée globale}, où l'intensité de traitement vaut 1 par construction et ne dérive pas.

\textbf{Ce qui ne sort pas} : les résultats v1 obtenus en portée partielle restent acquis et publiés ; ils ne sont simplement pas prolongés. La portée \code{new} (les naissances postérieures à t\ensuremath{_0} reçoivent la nouvelle technologie) reste disponible et n'est pas concernée : ce n'est pas un traitement partiel de la population vivante.

\subsection*{3.4 Purger les clefs mortes du carnet --- \textbf{exigence, pas option}}

Consigne explicite de l'utilisateur, 21 août 2026 : « \emph{Il est essentiel de purger les clefs mortes. Le run est en \ensuremath{\lambda}T\ensuremath{^2} en ce moment.} »

\textbf{Le diagnostic.} \code{LoanBook.by\_borrower} est un \code{defaultdict[int, set[int]]} (\code{m4\_3live/model.py:307}), et \code{remove()} retire un prêt par \code{.discard()} (\code{model.py:349}) : \textbf{la clef reste}, avec un ensemble vide. Une entité qui meurt voit tous ses prêts retirés (\code{\_fail\_one}, \code{model.py:608-616}) et laisse donc derrière elle une entrée vide, définitive. La phase de service des intérêts itère sur \textbf{toutes} les clefs --- \code{for borrower in list(book.by\_borrower.keys())} (\code{model.py:901}) --- et fait \code{continue} sur les vides (\code{model.py:903-904}).

\textbf{Le coût, correctement énoncé.} Le nombre de clefs croît comme le nombre d'entités jamais créées, soit \ensuremath{\approx} \ensuremath{\lambda}\ensuremath{\cdot}t. Le coût \emph{par pas} de la phase d'intérêts est donc en \ensuremath{\lambda}\ensuremath{\cdot}t, et le coût \textbf{total d'un run} est la somme de \ensuremath{\lambda}\ensuremath{\cdot}t sur les T pas, soit \textbf{\ensuremath{\lambda}T\ensuremath{^2}/2 : quadratique en T}. Mesuré en v1 : +20 \% de coût par pas entre t = 210 et t = 1810, à prêts actifs et population constants, avec 98 \% des clefs vides. À T = 4000 et \ensuremath{\lambda} = 30, le carnet porte \ensuremath{\approx} 120 000 clefs pour \ensuremath{\approx} 1130 entités vivantes. C'est le plafond qui empêche aujourd'hui les horizons longs.

\textbf{Ce qu'il faut faire.} Supprimer l'entrée de \code{by\_borrower}, \code{by\_lender}, \code{due}, \code{claims} et \code{debts} d'une entité \textbf{au moment de sa mort}, dans \code{\_fail\_one}, après le retrait de ses prêts. Reconstruire les listes de clefs une fois par pas devient alors proportionnel à la population vivante, et le run redevient linéaire en T.

\textbf{Et la parité tient --- c'est vérifiable, pas à parier.} Une note de v1 affirmait que purger « changerait l'ordre de sommation et ferait perdre la parité ». \textbf{Cette inférence est fausse} pour le cas qui nous occupe, pour trois raisons qui doivent toutes être vraies :

\begin{enumerate}
\item le corps de boucle est un \textbf{no-op} pour un ensemble vide --- le \code{continue} précède la lecture de \code{book.due[borrower]} (\code{model.py:903-905}), donc aucune opération flottante n'a lieu ;
\item supprimer une clef d'un \code{dict} CPython \textbf{ne réordonne pas} les clefs restantes ; la séquence des emprunteuses non vides est identique ;
\item une entité morte est tuée (\code{population.kill}) et \textbf{ne peut plus jamais emprunter}, donc sa clef ne sera pas réinsérée.
\end{enumerate}

\textbf{Le piège à ne pas commettre} : purger les clefs vides des entités \textbf{vivantes} (celles dont tous les prêts se sont simplement clos) \textbf{casserait} la parité --- un \code{defaultdict} réinsère la clef en \textbf{fin} d'ordre au prochain emprunt, ce qui change la séquence d'itération et donc l'ordre de sommation. Ne purger qu'à la mort.

\textbf{Porte de sortie} : parité bit à bit reconduite (lot A), et courbe de coût par pas mesurée avant/après sur un run de 4000 pas, pour montrer que la pente s'aplatit.

\subsection*{3.5 Une seule borne de transfert : \code{optimum}}

Consigne explicite de l'utilisateur, 21 août 2026 : « \emph{Je ne veux plus que \code{transfer\_cap="optimum"}.} »

La valeur \code{equalization} (plafond du transfert à (K\_\ensuremath{\ell} \ensuremath{-} K\_b)/2) est retirée. À supprimer, avec les références v1 :

\begin{center}
\begin{tabularx}{\linewidth}{@{}XX@{}}
\toprule
\textbf{À supprimer} & \textbf{Où (v1)} \\
\midrule
tuple \code{TRANSFER\_CAPS} et sa validation & \code{m4\_3live/model.py:69, 167-168} \\
champ \code{transfer\_cap} de \code{Config} & \code{m4\_3live/model.py:144} \\
\code{cap\_at\_equalization} et la branche de plafonnement & \code{m4\_3live/model.py:506, 531-536} \\
compteur \code{mkt\_capped} (devient mort) & \code{\_run\_market}, colonne de \code{series.csv} \\
option \code{-{}-transfer-cap} du CLI & \code{driver/headless.py:214} \\
champ exposé par l'API locale & \code{web/router.py:55} \\
pilote de plafond et sa figure \code{cap\_pilot.png} & \code{scripts/campaign.py:166}, \code{scripts/conception\_evidence.py:309, 348, 723} \\
contrôle d'intégrité « transfer\_cap inattendu » & \code{scripts/analyse.py:249-256} \\
\bottomrule
\end{tabularx}
\end{center}

\textbf{C'est du code mort, et c'est prouvé} : \code{mkt\_capped} cumulé vaut \textbf{exactement 0 sur les 203 runs} du dépôt. La suppression ne peut donc changer aucun nombre, et \textbf{la parité doit être reconduite telle quelle}. C'est ce qui la distingue du §3.1. À faire dans le lot A, avant tout changement de comportement, pour que la parité qui suit porte sur un moteur déjà simplifié.

Le pilote de plafond du 16 août reste publié dans le rapport de conception de v1 : il documente \emph{pourquoi} le plafond est sans objet (les entités dopées deviennent en quelques dizaines de pas le côté riche de presque toutes leurs paires, l'optimum voudrait alors leur envoyer du capital, et l'ancienne règle de sens l'interdisait). Ce résultat justifie la suppression ; il ne disparaît pas avec elle. \textbf{Noter la boucle} : c'est précisément le §3.1 qui rend ce raisonnement caduc, puisque le sens n'est plus interdit. Si le plafond redevenait pertinent sous sens libre, le dire dans le rapport de conception plutôt que de le réintroduire en silence.

\bigskip

\setcounter{section}{3}
\section{Instrumentation à ajouter au moteur forké}

\subsection*{4.1 Tension native, et l'effectif producteur exact}

\textbf{Note [17], la demande la plus explicite de l'utilisateur} --- « TRÈS IMPORTANT ». v1 l'a satisfaite \emph{a posteriori}, par un dérivé des séries (\code{m4\_3live\_credit\_soc/scripts/tension.py}), parce que le moteur était gelé. v2 n'a pas cette contrainte : \textbf{la tension doit être une colonne native}.

Ce qui manque dans v1 et qu'il faut enregistrer : par technologie et par pas, \code{n\_prod} (l'effectif qui a \textbf{produit} au pas) et \code{K\_prod} (le capital au moment de produire). Aujourd'hui \code{tech\_series} donne l'effectif de \emph{fin} de pas, alors que la production a été calculée \textbf{avant} les morts du pas ; \code{scripts/tension.py} reconstruit l'effectif producteur par \code{n + deaths}, ce qui est exact tant qu'une seule technologie vit et approché sinon (imputation au prorata). La colonne \code{basis} de \code{tension\_agg.csv} documente laquelle des reconstructions a servi --- cette colonne doit disparaître en v2, parce que la mesure devient exacte.

Sorties attendues, pour \textbf{chaque} run et sans intervention manuelle : \code{tension.csv} (par pas), \code{tension\_agg.csv} (agrégé) et une figure de tension dans \code{figures/}. Reprendre la structure de \code{m4\_3live\_credit\_soc/scripts/tension\_figures.py}, qui est branchée dans \code{driver/headless.py:write\_outputs()}.

\subsection*{4.2 Amplitude exacte, enregistrée et non reconstruite}

\textbf{Note [16].} Après le pas d'intervention, \code{population.prod[i] = A\ensuremath{'}\ensuremath{\cdot}K\_i\textasciicircum{}\{\ensuremath{\gamma}\ensuremath{'}\}} est en mémoire, donc \code{K\_i = (prod\_i/A\ensuremath{'})\textasciicircum{}\{1/\ensuremath{\gamma}\ensuremath{'}\}} exactement (aller-retour flottant propre), d'où la production contrefactuelle \code{A\ensuremath{\cdot}K\_i\textasciicircum{}\ensuremath{\gamma}} et donc l'amplitude \code{m} et la part traitée \code{p} \textbf{sans aucune inversion}. v1 l'a fait en post-traitement (\code{scripts/exact\_amplitude.py}) ; v2 doit l'enregistrer au moment de l'intervention.

Ce que ça a rapporté en v1, à reproduire comme contrôle : m = 1,500000 / 1,250000 / 1,946184 (le troisième est un levier sur \ensuremath{\gamma}, dont l'amplitude vaut \code{K\textasciicircum{}\ensuremath{\Delta}\ensuremath{\gamma}} et \textbf{doit être mesurée}, pas imposée) ; et \code{E(h=1) = 1} à 1,0\ensuremath{\cdot}10\ensuremath{^-}\ensuremath{^1}\ensuremath{^5}.

\subsection*{4.3 Prédiction naïve du levier \ensuremath{\gamma}}

\textbf{Note [17].} Pour tout levier sur \ensuremath{\gamma}, enregistrer côte à côte l'amplitude mesurée et la prédiction naïve qu'on aurait faite en supposant K = 1 (c'est-à-dire amplitude = 1). L'écart est le contenu informatif : en v1 il valait 1,946 contre 1,000. C'est ce qui montre qu'un levier sur \ensuremath{\gamma} n'est pas un cadran de puissance.

\subsection*{4.4 Diagnostics du régime nouveau (§3.1)}

\begin{itemize}
\item \code{mkt\_blocked\_dir} \textbf{contrefactuel} (ce qu'il aurait valu sous l'ancienne règle) ;
\item part du capital total détenue par les entités en \textbf{position nette créancière} ;
\item corrélation entre le rendement marginal d'une entité et sa position nette.
\end{itemize}

\subsection*{4.5 Statistique : ce que « x \ensuremath{\sigma} » veut dire}

\textbf{Note [19].} Avec n graines, les écarts rapportés sont des \textbf{Student à n\ensuremath{-}1 degrés de liberté}, pas des normales. À 5 graines : seuil 5 \% à \textbf{2,776} et non 1,96 ; seuil 1 \% à \textbf{4,604} et non 2,58. v1 a vérifié qu'aucun verdict ne basculait, mais une marge était mince (p = 3,3 \%).

\textbf{Recommandation} : passer à \textbf{12 graines} pour la campagne de verdict, ce qui ramène le seuil à 2,201 et divise l'erreur-type par 1,55. Coût : \ensuremath{\times}2,4 sur la campagne. À confirmer contre le budget (§12).

\bigskip

\setcounter{section}{4}
\section{La question théorique : qu'est-ce qui détermine la rotation du crédit ?}

C'est le programme scientifique de v2, et il est ouvert.

\textbf{Ce que v1 a établi.} Sur 109 runs couvrant quatre leviers indépendants (K0 seul, \ensuremath{\delta} seul, A seul, A et K0 ensemble), \code{morts/population/pas \ensuremath{\propto} (loan\_volume/K\_tot)\textasciicircum{}1,337} avec R\ensuremath{^2} = 0,9982, écart médian 0,41 \%, et un biais par levier compris entre \ensuremath{-}1,7 \% et +0,7 \%. Aucune autre variable testée n'aligne les quatre leviers --- la tension échoue sur \ensuremath{\delta}.

\textbf{Ce qui manque.} La rotation est \textbf{endogène} : on la mesure, on ne la choisit pas. Tant qu'on ne sait pas la prédire à partir des paramètres, on a une régularité descriptive et non une loi.

\textbf{Piste, non imposée.} La rotation est le produit de trois facteurs : le nombre de paires tirées par pas \code{\ensuremath{\eta}(N) = \ensuremath{\rho}N} (\code{model.py:497}), la probabilité qu'une paire traite, et le transfert moyen rapporté au capital. Les deux premiers sont combinatoires et se calculent ; le troisième est \code{E|\ensuremath{\delta}*|/K}, dont l'institution donne la forme exacte. Une prédiction analytique est donc plausible. \textbf{Le §3.1 change les deux derniers facteurs} : plus aucune paire n'est refusée pour cause de sens, donc la probabilité de traiter monte et le transfert moyen change de distribution. C'est une raison de plus de faire le §3.1 d'abord.

Les 203 runs de v1 suffisent pour commencer l'analyse sans lancer un seul calcul neuf.

\bigskip

\setcounter{section}{5}
\section{Les deux groupes de covariance : à énoncer ensemble}

v1 les a démontrés séparément. \textbf{v2 doit les présenter comme un seul résultat} dans son rapport de conception : ce sont deux groupes de symétrie qui réduisent le nombre de paramètres réellement libres du modèle.

\textbf{Covariance d'échelle (capital)} --- rapport v1 §7.1. Avec \code{A \ensuremath{\rightarrow} A\ensuremath{'}} et \code{K0 \ensuremath{\rightarrow} cK0}, \code{c = (A\ensuremath{'}/A)\textasciicircum{}\{1/(1\ensuremath{-}\ensuremath{\gamma})\}}, on obtient la \textbf{même simulation à l'échelle c près} : chaque phase du pas est homogène de degré 1 en capital, et le taux marginal \code{\ensuremath{\gamma}AK\textasciicircum{}\{\ensuremath{\gamma}\ensuremath{-}1\}} est \textbf{sans dimension}, donc invariant. Vérifié à 5\ensuremath{\cdot}10\ensuremath{^-}\ensuremath{^1}\ensuremath{^5}, population rigoureusement identique. Ce qui casse l'exactitude : trois constantes dimensionnées non rééchelonnées (\code{MIN\_LOAN}, \code{ZERO\_TOL}, le plafond de transfert --- ce dernier disparaissant avec le §3.5).

\textbf{Covariance de pas de temps} --- rapport v1 §7.3. Avec un pas de recalage \code{s} (le nombre de pas anciens comprimés en un) : \code{\ensuremath{\lambda}\ensuremath{\rightarrow}s\ensuremath{\lambda}}, \code{\ensuremath{\delta}\ensuremath{\rightarrow}1\ensuremath{-}(1\ensuremath{-}\ensuremath{\delta})\textasciicircum{}s}, \code{\ensuremath{\sigma}\ensuremath{\rightarrow}\ensuremath{\sigma}\ensuremath{\sqrt{\;}}s}, \textbf{\code{A\ensuremath{\rightarrow}sA}}, \textbf{\code{\ensuremath{\rho}\ensuremath{\rightarrow}s\ensuremath{\rho}}}. Les trois premières substitutions composent \textbf{exactement} (identité algébrique ; somme de deux log-normales dérive comprise ; somme de deux Poisson). Les deux dernières sont les \textbf{flux par pas}, que l'énoncé initial oubliait.

\textbf{Classification à porter dans le \code{model\_summary} de v2} :

\begin{center}
\begin{tabularx}{\linewidth}{@{}XX@{}}
\toprule
\textbf{Grandeurs \textbf{par pas} (se recalent avec l'unité de temps)} & \textbf{Grandeurs \textbf{de forme} (n'en dépendent pas)} \\
\midrule
\ensuremath{\delta}, \ensuremath{\sigma}, \ensuremath{\lambda}, A, \ensuremath{\rho} & \ensuremath{\gamma}, K0, règle de transfert, règle de taux \\
\bottomrule
\end{tabularx}
\end{center}

Cette colonne évite de discuter un « effet de \ensuremath{\delta} » qui n'est qu'un changement d'unité de temps.

\textbf{Une limite mesurée, à ne pas réouvrir naïvement.} Le résidu du recalage (6 à 17 \%) ne vient \textbf{pas} de la discrétisation en \ensuremath{\delta} --- c'était l'hypothèse de départ, elle a été \textbf{réfutée}. Le test à \code{s = 4} ne tranche pas (les deux termes croissent en s\ensuremath{-}1) ; celui à \code{\ensuremath{\delta} = 0,002} tranche : le terme de K\_aut est divisé par 5 et le résidu ne bouge pas. La source est la \textbf{phase de marché, qui ne se compose pas} --- deux rondes de \ensuremath{\rho}N appariements ne valent pas une ronde de 2\ensuremath{\rho}N, la seconde voyant l'état laissé par la première. C'est ce qui distingue ce modèle d'une équation différentielle, et c'est la vraie question ouverte : \emph{que font s rondes successives qu'une ronde s fois plus grosse ne fait pas ?}

\bigskip

\setcounter{section}{6}
\section{Protocole expérimental}

\textbf{Reprendre le protocole v1 tel quel} (rapport §3), sauf sur les points ci-dessous. Il est éprouvé : amorçage partagé de 2000 pas, snapshots à t\ensuremath{_0} = 2000, bras branchés sur les mêmes snapshots, appariement par graine, fenêtre d'observation sur les 1000 ou 2000 derniers pas.

\textbf{Modifications pour v2} :

\begin{enumerate}
\item \textbf{Portée globale uniquement} (§3.3). Plus de bras \code{frac\_*}, plus de bras \code{null}. Le contrôle apparié redevient un simple bras inerte.
\item \textbf{12 graines} pour la campagne de verdict (§4.5), à confirmer.
\item \textbf{Comparaison v1/v2 appariée.} Le chantier du §3.1 exige une campagne A/B où le \emph{seul} changement est la règle de sens : mêmes graines, mêmes snapshots d'amorçage, même protocole. Attention : le fork v2 doit pouvoir \textbf{rejouer l'ancienne règle} (drapeau de configuration, \code{loan\_direction \ensuremath{\in} \{"richest\_lends", "free"\}}), sinon la comparaison n'est pas appariée mais inter-lignée.
\item \textbf{Contrôle de stationnarité obligatoire} avant de lire un niveau : rapport de la moyenne du dernier quart à celle du quart précédent, sur \code{K\_tot}, \code{pop} et \code{prod\_tot}. v1 accepte 0,993 à 1,003. Un bras qui n'est pas stationnaire donne un transitoire, pas un niveau --- le dire.
\end{enumerate}

\bigskip

\setcounter{section}{7}
\section{Séquencement en lots}

Chaque lot a une \textbf{porte de sortie} : tant qu'elle n'est pas franchie, le lot suivant ne commence pas. Les lots qui changent le modèle doivent être \textbf{séparés dans l'historique git} et chacun accompagné de sa campagne appariée. Ne jamais les empiler avant mesure --- c'est la seule manière de garder un verdict attribuable.

\begin{center}
\begin{tabularx}{\linewidth}{@{}lXXl@{}}
\toprule
\textbf{Lot} & \textbf{Contenu} & \textbf{Porte de sortie} & \textbf{Coût estimé} \\
\midrule
\textbf{A} & Fork du dossier ; \textbf{§3.4} purge des clefs mortes ; \textbf{§3.5} suppression de \code{equalization} ; parité sur le fork nu & \code{test\_parity\_m4\_3.py -{}-full} : \textbf{écart nul exigé} --- les deux changements sont neutres sur le calcul (§3.4, §3.5). \textbf{Et} courbe de coût par pas mesurée avant/après sur 4000 pas : la pente doit s'aplatir & 3 h + 2200 s de calcul \\
\textbf{B} & \textbf{§3.1} sens du prêt libre, nomenclature \code{K\_a}/\code{K\_b}, compteur contrefactuel, diagnostics §4.4 & parité reconduite \textbf{ou} écart expliqué ligne à ligne --- pas de tolérance de confort & 1 j \\
\textbf{C} & \textbf{§4.1--4.3} tension native, amplitude exacte, prédiction naïve \ensuremath{\gamma} & tests unitaires sur données synthétiques dont on connaît la réponse & 0,5 j \\
\textbf{D} & Campagne A/B appariée du §3.1 : ancienne règle vs sens libre, 12 graines & verdict sur le régime nouveau, avec sa statistique de Student & 2 h de calcul \\
\textbf{E} & \textbf{§3.2} \code{phase\_order}, campagne A/B appariée & taux de défaut mesuré \textbf{contre la prédiction écrite a priori} & 0,5 j + 2 h de calcul \\
\textbf{F} & \textbf{§5} analyse de la rotation du crédit sur les runs existants + ceux du lot D & soit une prédiction analytique testée, soit un constat d'échec documenté avec ce qui a été essayé & 1 j \\
\textbf{G} & \textbf{§6} énoncé unifié des deux covariances dans le rapport de conception & les deux groupes présentés ensemble, avec la classification par pas / de forme & 0,5 j \\
\textbf{H} & Rapports (§10), journal, traçabilité, import \code{simulation\_lab} & relecture annotée par l'utilisateur & 1,5 j \\
\bottomrule
\end{tabularx}
\end{center}

\emph{(Il n'y a pas de lot « portée persistante » : le cas partiel est hors périmètre, §3.3.)}

\bigskip

\setcounter{section}{8}
\section{Tests exigés}

Assertions Python simples, convention du dépôt (\textbf{pas de pytest}). Reprendre la suite de v1 (\code{m4\_3live\_credit\_soc/tests/}) et l'étendre.

\subsection*{Trois sémantiques de parité, à ne pas confondre}

C'est le point le plus facile à saboter par inattention.

\begin{itemize}
\item \textbf{Lot A --- parité obligatoire et bit-exacte.} Supprimer \code{equalization} retire du code qui n'a jamais été exécuté (\code{mkt\_capped} = 0 partout). Toute déviation est un bug de la suppression, pas un effet du changement. Commande : \code{python3 tests/test\_parity\_m4\_3.py -{}-full} --- 8000 pas, 26 colonnes, \textasciitilde{}1100 s, référence \code{m4\_3\_\_d1\_\_baseline\_\_seed0}.
\item \textbf{Lot B --- parité à retester, pas à supposer.} Le sens libre coïncide avec l'ancienne règle en régime homogène, mais l'ordre du couple change l'ordre de sommation et d'insertion dans le carnet. Si la parité tombe, \textbf{expliquer l'écart ligne à ligne} ; ne pas l'absorber dans une tolérance.
\item \textbf{Lot E --- parité conditionnelle.} Le défaut \code{phase\_order = "v1"} doit la préserver ; la branche \code{"deprec\_first"} ne la préserve pas, par construction, et c'est le but.
\end{itemize}

\subsection*{Autres tests}

\begin{itemize}
\item \textbf{Institution} : le régime homogène applique littéralement (K\_b \ensuremath{-} K\_a)/2 ; les régimes hétérogènes vérifiés contre la formule fermée et contre une résolution de Newton.
\item \textbf{Sens du prêt} (nouveau) : sur une paire construite à la main où l'optimum exige que la pauvre cède, vérifier que le prêt part bien dans ce sens, que le carnet enregistre la bonne débitrice, et que le compteur contrefactuel s'incrémente.
\item \textbf{Rejeu/déterminisme} : un journal issu d'une \textbf{session en direct réelle} (interventions soumises de façon asynchrone), rejoué en mode headless, reproduit une trajectoire strictement identique, même \code{t} effectif compris.
\item \textbf{Invariant de pause} : N pas avec pauses = N pas sans pause.
\item \textbf{Portées} : \code{all} change bien toutes les entités vivantes au pas suivant ; \code{new} ne change rien aux entités déjà nées.
\item \textbf{Tension} (nouveau) : sur un état synthétique à une seule technologie dont on connaît \code{n}, \code{K} et \code{A}, vérifier que \code{K\_eq} et \code{T} valent exactement la formule.
\item \textbf{Covariance de pas de temps} (§12, question ouverte) : candidat naturel pour un test de non-régression --- mais son résidu est de 6 à 17 \%, donc le seuil devrait être calibré sur le résidu mesuré, pas sur zéro.
\end{itemize}

\bigskip

\setcounter{section}{9}
\section{Livrables et conventions de rédaction}

L'utilisateur a explicitement demandé que v2 reproduise la manière de v1 : \textbf{un rapport scientifique, illustré, détaillé et autonome ; un rapport de conception ; un journal de travail.} Ce qui suit n'est donc pas indicatif.

\subsection*{Livrables}

\begin{itemize}
\item \textbf{Code complet et fonctionnel} dans \code{m4\_3live\_v2\_credit\_soc/} --- un système qui tourne réellement, pas un squelette. Démarrer le serveur, piloter une session en direct dans un vrai navigateur, vérifier visuellement les graphiques et le panneau de contrôle avant de déclarer la tâche terminée.
\item \textbf{Rapport de conception}, LaTeX compilé en PDF (\code{report/conception\_m4\_3live\_v2.pdf}) : toutes les décisions d'architecture justifiées avec citations \code{fichier.py:ligne} --- en particulier le sens libre et sa nomenclature (§3.1), le statut de la parité après chaque lot, la purge des clefs mortes (§3.4) et la suppression de \code{equalization} (§3.5), l'instrumentation native (§4), l'énoncé unifié des deux covariances (§6), et les décisions laissées ouvertes (§12) avec la façon dont tu les as tranchées.
\item \textbf{Rapport de résultats}, LaTeX compilé en PDF (\code{report/rapport\_final.pdf}) : protocole, résultats, figures, verdict, limites. Même discipline que \code{m4\_3live\_credit\_soc/report/rapport\_final.tex}, qui est le modèle à imiter.
\item \textbf{Compilation PDF vérifiée} --- le \code{.tex} seul ne suffit pas. Trois passes \code{pdflatex} après suppression des \code{.aux}/\code{.toc}, aucun \code{!} dans le log, aucune référence non définie.
\item \textbf{\code{README.md}} et \textbf{\code{JOURNAL.md}} dans \code{m4\_3live\_v2\_credit\_soc/}. Le journal est tenu \textbf{au fil de l'eau}, pas reconstitué à la fin : chaque session y consigne ses chiffres clés, ses décisions, et les inférences qu'elle a dû retirer.
\end{itemize}

\subsection*{Conventions de rédaction --- règles permanentes du dépôt}

\textbf{Autonomie.} Un rapport doit être compréhensible par un lecteur tiers compétent dans le domaine général, \textbf{sans accès aux données, au code, ni aux conversations qui l'ont précédé}.

\begin{itemize}
\item Toute variable, tout paramètre, toute notation est \textbf{défini explicitement à sa première apparition} --- même si la définition semble évidente dans le contexte de travail. Jamais de notation avant sa définition.
\item Tout terme technique non standard est expliqué dans le corps du texte.
\item Toute méthode d'analyse est décrite avec assez de détail pour qu'un lecteur puisse la réimplémenter ou en évaluer la validité ; les hypothèses implicites sont rendues explicites.
\item Toute valeur numérique et toute figure a une \textbf{source identifiable} dans le texte (section, tableau, fichier de données). Les chiffres ne tombent pas du ciel.
\item Une conclusion est \textbf{proportionnée aux données} : distinguer ce qui est robuste (confirmé sur plusieurs graines, paramètres, horizons) de ce qui est fragile ou exploratoire.
\end{itemize}

\textbf{Marquage typographique obligatoire.} v1 emploie trois macros LaTeX qui rendent la règle des 95 \% visible à la lecture : \code{\textbackslash\{\}fait\{\}} (fait observé), \code{\textbackslash\{\}inference\{\}} (inférence), \code{\textbackslash\{\}incertitude\{\}} (incertitude ou réserve). Elles s'affichent en clair dans le PDF. \textbf{Les reprendre.} C'est une grande part de ce que l'utilisateur entend par « rapport scientifique » : on doit pouvoir lire un paragraphe et savoir immédiatement s'il rapporte une mesure ou une interprétation.

\textbf{Cinq règles tirées des annotations de l'utilisateur sur le rapport v1} :

\begin{enumerate}
\item \textbf{Tout tableau de plus de quatre colonnes est d'abord une figure.} Si l'information est comparative, elle se voit ; le tableau vient après, pour les chiffres exacts.
\item \textbf{Aucune notation avant sa définition}, sans exception.
\item \textbf{Un calcul cité est un calcul déroulé.} Ne pas écrire « soit \ensuremath{\times}1,80 » : écrire l'opération.
\item \textbf{Un paragraphe qu'on ne peut pas paraphraser est à réécrire, pas à défendre.}
\item \textbf{Jamais de « meilleure technologie ».} L'ordre des rendements marginaux dépend du capital : il n'existe pas de technologie meilleure dans l'absolu. Écrire « plus grand rendement marginal à ce capital ».
\end{enumerate}

\textbf{Traçabilité par hash --- règle permanente.} Toute simulation mentionnée dans un rapport (figure, tableau ou texte) doit être identifiable par son \code{run\_id}. Chaque rapport se termine par une \textbf{annexe de traçabilité} donnant, pour chaque run : \code{run\_id}, l'endroit où il apparaît, son rôle, ses paramètres principaux, sa graine, son statut, son chemin \code{simulation\_lab}. v1 génère cette annexe par script (\code{scripts/make\_traceability.py}) et vérifie qu'aucune ligne n'est sans rôle --- reprendre le mécanisme. \textbf{Dans cette annexe, renvoyer aux sections, pas aux numéros de figure} : v1 s'est fait piéger, ses renvois dataient d'une version où le rapport avait deux figures de moins.

\textbf{Travail scientifique --- règles permanentes.}

\begin{itemize}
\item \textbf{Données} : toute donnée qui a informé une décision est sauvegardée dans un fichier accessible (JSON ou CSV, sous \code{results/}), pas seulement dans un \code{print} ou un log.
\item \textbf{Raisonnement} : toute inférence non triviale est explicitée avec les chiffres qui la justifient. Une conclusion sans source chiffrée est invalide.
\item \textbf{Reproductibilité} : tout calcul, toute analyse, toute figure est accompagné du script qui permet de le reproduire. Un résultat non reproductible est un résultat non publié. Le rapport se termine par une section « Reproduire » listant les commandes dans l'ordre, avec leur coût mesuré.
\item \textbf{Journal} : consigné immédiatement, avec les chiffres clés.
\end{itemize}

\bigskip

\setcounter{section}{10}
\section{Arborescence attendue}

\begin{quote}\footnotesize
\begin{verbatim}
m4_3live_v2_credit_soc/
+-- ROADMAP.md                       (feuille de route, deja presente)
+-- notes/                           (28 notes extraites du PDF v1, deja presentes)
+-- prompts/PROMPT_M4_3LIVE_V2.md    (ce document)
+-- m4_3live_v2/                     (paquet moteur FORKE, S0 -- nom a confirmer S12)
|   +-- model.py                     (sens libre S3.1, phase_order S3.2, instrumentation S4)
|   +-- kernel.py                    (repris tel quel de v1)
|   +-- live.py                      (session, boucle, journal)
+-- driver/                          (pilote headless)
+-- web/                             (IHM, ou extension de simulation_lab/live/)
+-- tests/
|   +-- test_parity_m4_3.py          (S9, trois semantiques)
|   +-- test_institution.py
|   +-- test_loan_direction.py       (S9, NOUVEAU)
|   +-- test_replay.py
|   +-- test_scopes.py
|   +-- test_tension.py              (S9, NOUVEAU)
+-- scripts/                         (campagnes et analyses)
+-- results/                         (runs, snapshots, analysis/)
+-- report/
|   +-- conception_m4_3live_v2.tex/.pdf
|   +-- rapport_final.tex/.pdf
|   +-- figures/
|   +-- tables/
+-- README.md
+-- JOURNAL.md
\end{verbatim}
\end{quote}

\bigskip

\setcounter{section}{11}
\section{Décisions laissées ouvertes, à trancher et documenter toi-même}

\begin{enumerate}
\item \textbf{Nom du paquet forké} --- \code{m4\_3live\_v2} par défaut ; changer si gênant.
\item \textbf{Le taux d'intérêt quand la créancière est la plus fragile} (§3.1) --- la règle de taux actuelle (moyenne géométrique des rendements marginaux) a-t-elle encore un sens si celle qui prête est la plus pauvre ? À trancher avant le lot D, et à documenter quelle que soit la réponse.
\item \textbf{Nombre de graines de verdict} (§4.5) --- 12 recommandé, à confirmer contre le budget de calcul.
\item \textbf{Constantes dimensionnées} (§6) --- rééchelonner \code{MIN\_LOAN} et \code{ZERO\_TOL} en unités de capital rendrait la covariance d'échelle exacte, mais casse la comparabilité directe avec toutes les campagnes antérieures.
\item \textbf{La covariance de pas de temps devient-elle un test de non-régression ?} (§6, §9) --- au même titre que la parité, ou reste-t-elle une vérification ponctuelle du rapport ? Attention : son résidu est de 6 à 17 \%, donc un seuil naïf à zéro échouerait toujours.
\end{enumerate}

\bigskip

\setcounter{section}{12}
\section{Pièges de mesure hérités de v1}

Ceux qui sont \textbf{toujours actifs} :

\begin{itemize}
\item \textbf{L'amplitude d'un levier sur \ensuremath{\gamma} vaut \code{K\textasciicircum{}\ensuremath{\Delta}\ensuremath{\gamma}}} (\ensuremath{\approx} 1,95 en v1, pas 1,2) : elle dépend du capital et doit être \textbf{mesurée à l'horizon 1}, jamais imposée.
\item \textbf{Le coût d'un run est QUADRATIQUE en T}, et pas seulement croissant avec la taille du système : le carnet accumule une clef morte par entité jamais créée, la phase d'intérêts les parcourt toutes, donc le coût par pas est en \ensuremath{\lambda}\ensuremath{\cdot}t et le coût total en \textbf{\ensuremath{\lambda}T\ensuremath{^2}/2}. Mesuré en v1 : +20 \% de coût par pas entre t = 210 et t = 1810 à prêts actifs et population constants, 98 \% des clefs étant vides. \textbf{C'est corrigé en v2 (§3.4), ce n'est pas une fatalité à contourner.}
\item \textbf{Le choc multiplicatif ne biaise pas une comparaison appariée.} Il frappe entre l'action et la production, mais il est \textbf{commun aux deux bras} : mêmes vivantes, même ordre, même état du générateur. Il ne s'annule pas terme à terme --- il est identique des deux côtés. Corollaire à préserver en v2 : \emph{tant qu'une intervention ne touche pas au capital, tout ce qui est en amont de la production se simplifie dans un rapport apparié.} C'est ce qui rend l'appariement par graine puissant.
\item \textbf{Pour mesurer l'ordre de convergence d'une table, l'échantillon doit balayer la variable indexée}, sinon on ne sonde qu'un point.
\end{itemize}

Ceux qui deviennent \textbf{dormants} avec le retrait du cas partiel (§3.3), à connaître mais pas à outiller --- sauf si une portée partielle réapparaît un jour :

\begin{itemize}
\item un tirage \code{fraction} consomme le générateur, donc sa référence appariée est un bras \code{null} qui tire la même cohorte sans rien appliquer, jamais un contrôle inerte ;
\item la part de production d'une cohorte traitée se mesure \textbf{après} application du levier, et doit être inversée en part ex ante par \code{p = s/(m \ensuremath{-} (m\ensuremath{-}1)s)} avant toute élasticité, sinon on la sous-estime de 10 \% ;
\item la portée \code{fraction} ne fixe \textbf{pas} l'intensité de traitement dans ce modèle : elle décroît vers zéro en \textasciitilde{}300 pas.
\end{itemize}

\bigskip

\setcounter{section}{13}
\section{Quatre leçons d'analyse, générales}

Elles ont coûté du temps réel en v1 et elles dépassent ce modèle.

\begin{enumerate}
\item \textbf{Un garde-fou d'ajustement doit tester l'ÉTENDUE, pas le nombre de valeurs distinctes.} v1 traçait des régressions sur neuf runs couvrant une plage de 0,6 \% --- du bruit de graine --- parce que le garde-fou comptait les valeurs distinctes.
\item \textbf{Un ajustement groupé sur plusieurs familles peut n'être que la droite qui joint les lignes de base.} v1 avait un R\ensuremath{^2} de 0,72 sur 127 runs ; la droite passant par les \textbf{trois bases seules} donne l'exposant 0,410 contre 0,396 pour l'ajustement complet. Il ne mesurait aucune réponse. \textbf{Contrôle systématique} : ajuster sur les seules bases et comparer.
\item \textbf{Un observable peut être ininformatif précisément parce que deux effets se compensent.} Sous le recalage temporel brut, \code{K\_tot} ne bouge que de 1 \% --- non parce que le système est conservé, mais parce que chaque entité est deux fois plus petite et qu'il y en a presque deux fois plus.
\item \textbf{Une identité qui est une réécriture des définitions ne doit jamais être présentée comme une confirmation empirique.} L'identité \code{dlnT/dlnA = 1/(1\ensuremath{-}\ensuremath{\gamma}) \ensuremath{-} (1/\ensuremath{\gamma})(\ensuremath{\varepsilon}\_prod \ensuremath{-} \ensuremath{\eta}\_pop \ensuremath{-} 1)} ne pouvait pas être fausse ; son résidu de 10\ensuremath{^-}\ensuremath{^3} mesure la propreté de l'aller-retour numérique, rien d'autre. Le dire dans le texte.
\end{enumerate}

À quoi s'ajoute une \textbf{habitude à prendre} : quand un test est possible, \textbf{écrire la prédiction avant de lancer les runs}, et rapporter ensuite ce que la prédiction a et n'a pas attrapé. v1 l'a fait deux fois ; les deux fois la prédiction était partiellement juste, et le dire valait mieux que de ne rapporter que la mesure.

\bigskip

\setcounter{section}{14}
\section{Consignes de méthode de travail}

\begin{itemize}
\item \textbf{Livrer un système complet, pas des ébauches.} Chaque pièce du §11 doit fonctionner de bout en bout avant que la tâche soit déclarée terminée --- pas de squelette avec des TODO, pas de fonctionnalité à moitié câblée.
\item \textbf{Calibrer la longueur des livrables écrits sur leur substance.} Les deux rapports doivent couvrir ce que ce prompt demande, sans sections de remplissage, sans résumé redondant en fin de document, sans paragraphe qui répète ce qui vient d'être dit autrement.
\item \textbf{Paralléliser les calculs par lot.} La machine a 8 cœurs, dont \textbf{7 utilisables} (autorisation explicite de l'utilisateur, 21 août 2026) : \textbf{7 processus au maximum}, un par lot de runs indépendants. Les coûts ci-dessous ont été mesurés à 6 processus ; les rapporter à 7 en le disant. Ordres de grandeur mesurés en v1 : amorçage 5 graines \ensuremath{\times} 2000 pas = 220 s ; 9 bras \ensuremath{\times} 5 graines \ensuremath{\times} 2000 pas = 1865 s ; ablation 25 runs = 1287 s ; parité 8000 pas = 1077 s. Attention : un run à \ensuremath{\gamma} = 0,4 coûte \textbf{4 fois} un run à \ensuremath{\gamma} = 0,5, parce que la population double et que le marché et le service doublent avec elle.
\item \textbf{Périmètre : livrer ce qui est demandé, à l'échelle demandée.} Prendre seul les décisions d'implémentation de routine (y compris celles du §12) ; ne pas élargir le mandat au-delà du §1 ; si une lecture différente de ce prompt changerait matériellement le travail à produire, le signaler en une phrase et continuer avec la lecture la plus proche du texte.
\item \textbf{Signaler ce qui est retiré.} Si une inférence écrite plus tôt se révèle fausse, la retirer explicitement dans le journal avec ce qui l'a réfutée. v1 l'a fait trois fois ; c'est une partie du travail, pas un aveu.
\end{itemize}

\subsection*{Sous-agents : subdiviser intelligemment, pas systématiquement}

\begin{itemize}
\item \textbf{Réserver ton propre raisonnement aux décisions qui déterminent la validité scientifique du travail} : la sémantique du sens libre et son effet sur le carnet (§3.1), les trois sémantiques de parité (§9), le protocole et l'écriture des verdicts (§7), l'analyse de la rotation (§5), et la relecture finale de cohérence entre code et rapports.
\item \textbf{Déléguer les missions simples, bien délimitées et vérifiables isolément} à des sous-agents sur un modèle plus léger : écrire un test unitaire donné une spécification précise, lancer et collecter une graine d'une campagne dont le protocole est gelé, rédiger une section descriptive à partir de décisions déjà prises et citées, exécuter une non-régression sur \code{simulation\_lab} et rapporter le résultat brut.
\item Donner à chaque sous-agent un mandat \textbf{autonome et complet} --- quel fichier, quel contrat d'entrée/sortie, quel test de validation. Un sous-agent frais ne voit pas ce prompt : lui fournir directement les extraits pertinents (numéros de ligne, formules, sémantique) plutôt que d'y renvoyer par référence.
\item Ne pas déléguer la vérification de ton propre travail à un sous-agent en plus des tests déjà exigés (§9).
\item Ne pas ouvrir plus de sous-agents que la structure du travail n'en justifie.
\end{itemize}


\end{document}
